\subsection{Ambiente, criteri e raccolta delle evidenze}
Ambito: UC-02, RF-10, RF-11, RF-17, RNF-02 e feedback UI di RNF-05.
Sono esclusi da questo piano export, immagini e archivio, coperti da
verifiche di integrazione separate. L'ambiente comprende .NET 10, xUnit
e PostgreSQL 16 dedicato ai test sulla porta 55432.
Il controllo dal browser richiede inoltre Angular e l'API del setup di prova.
Gli account di prova e gli ID sono predisposti dalle fixture. I casi
di rifiuto confrontano snapshot persistenti prima e dopo la richiesta;
account inattivo e issue duplicata sono condizioni separate.

\textbf{Criteri di ingresso:} disponibilita' di API, login e database isolato,
con data, riferimento dei sorgenti, versioni, ambiente e comando registrati.
L'assenza di un prerequisito comporta un esito bloccato.
La suite HTTP usa database distinti per caso e JWT ottenuti dal login reale.
Una fixture puo' predisporre condizioni non ottenibili dalla UI, per esempio
un READONLY assegnatario, senza modificare i dati di sviluppo.

\textbf{Oracoli:} stato, date UTC, destinatario e notifiche persistenti.
Due chiamate separate a \texttt{UtcNow} non devono necessariamente produrre
date identiche. Stesso stato valido: date e notifiche restano invariate;
rifiuto: nessuna scrittura. La verifica associa alla risposta HTTP il
riscontro persistente e l'evidenza dell'esito. Il controllo del proprietario
si applica anche quando chi tenta la lettura e' ADMIN.

\textbf{Evidenze acquisite:} suite backend, coverage, persistenza e analisi
qualita' sono in \path{verifica-backend/} sotto Documentazione finale.
Il report frontend storico e' in \path{verifica-frontend/}; l'audit del
19 ottobre registra i nuovi esiti. RF-17 e' verificato in backend e UI.
Il CSV associa i casi ai metodi passati nel TRX e ai riscontri browser,
distinti dalle tre sessioni con utenti descritte nella valutazione di usabilita'.

La tabella identifica quattordici casi, TP-01--TP-14. Stato e permessi sono in
\path{IssueStatusIntegrationTests.cs}; duplicati in
\path{IssueDuplicateValidationIntegrationTests.cs}; notifiche in
\path{NotificationsIntegrationTests.cs}. Riapertura/seconda risoluzione,
account disattivato e ripetizione del PATCH read hanno esiti dedicati
nel TRX. Per TP-14, RF-17 fissa l'oracolo: GET delle
proprie notifiche 200; PATCH read su notifica propria esistente 403, anche
se gia' letta, senza modificare \texttt{IsRead}, date o altre righe.
Le fixture includono per R notifiche lette e non lette. Il riscontro browser
verifica l'assenza del comando UI e di PATCH automatici durante la consultazione;
USER e ADMIN mantengono la lettura delle proprie notifiche (TP-13 e regressione
per ADMIN). Esiti e riferimenti sono nel CSV; il criterio RF-17 e' invariato.
